\bookchapter{제18장: 데이터 백업과 이전}

\chapterintro{이웃 한 분이 스마트폰을 물에 빠뜨려 꺼냈더니 이미 부팅이 안 됐습니다. 그녀가 가장 마음 아파한 것은 스마트폰 자체가 아니라, 3년간 백업하지 않은 사진——아이가 태어나서 걷기 시작한 전체 기록이었습니다. 이야기를 듣고 마음이 정말 편치 않았습니다, 이 일은 완전히 피할 수 있었기 때문입니다. 이 장에서는 진지하게 데이터에 보험을 들어보겠습니다.}

\booksection{하나, 먼저 파악: 스마트폰 안의 어떤 것을 잃어버리면 되돌릴 수 없나}

데이터마다 “생존 상태”가 다릅니다, 구분해야 합니다:

\begin{itemize}
  \item \textbf{스마트폰 로컬에 저장, 잃어버리면 진짜 없어지는 것}: 사진 비디오(클라우드 동기화 안 한 경우), 메모, 문자, 통화 기록, 다운로드한 파일.
  \item \textbf{앱 서버에 저장, 스마트폰 바꿔도 다시 로그인하면 있는 것}: 위챗 친구 목록, 타오바오 주문, 은행 앱의 거래 기록. 이것들은 해당 회사의 서버에 저장, 스마트폰에 있는 게 아님.
  \item \textbf{회색 지대, 특별한 조작이 필요한 것}: 위챗 채팅 기록(위챗 서버에 저장되지만, 본인이 서버에서 다운로드할 수 없음), 스마트폰 뱅킹 바인딩 기기 증명서, 2단계 인증 앱의 키(Google Authenticator 같은).
\end{itemize}

그래서 당신이 중점적으로 보호해야 할 것은, 첫 번째와 세 번째 부류입니다.

\booksection{둘, 전체 백업: 한 번에 모든 것 해결}

전체 백업은 스마트폰의 완전한 “스냅샷”을 찍는 것과 같습니다——현재 시스템 설정, 데스크톱 레이아웃, 앱 데이터, 사진, 연락처, 전부 패키징.

\booksubsection{방안 1: 컴퓨터 로컬 백업(가장 신뢰, 컴퓨터 있는 사람 권장)}

\begin{itemize}
  \item \textbf{iPhone}: 스마트폰을 데이터 케이블로 컴퓨터에 연결, iTunes(Windows 구버전) 또는 Finder(Mac) 열기, 기기 페이지에서 "지금 백업" 선택. "백업 암호화" 체크 주의, 그렇지 않으면 비밀번호와 건강 데이터가 백업에 포함되지 않음.
  \item \textbf{Android}: 각 브랜드마다 자체 컴퓨터용 소프트웨어, 예: 화웨이 폰 어시스턴트, 샤오미 어시스턴트, 삼성 Smart Switch. 컴퓨터 연결 후 안내에 따라 조작.
  \item 컴퓨터 백업의 장점은 클라우드 공간 살 필요 없음, 그리고 복원 속도가 클라우드 복원보다 훨씬 빠름. 단점은 정기적으로 케이블 연결 백업 기억 필요——대부분 이것을 못 함.
\end{itemize}

\booksubsection{방안 2: 클라우드 백업(가장 편리, 모든 사람 권장)}

\begin{itemize}
  \item \textbf{iPhone}: 설정 본인 이름 iCloud iCloud 클라우드 백업 켜기. 스마트폰이 전원 연결, Wi-Fi 연결, 잠금 화면 상태일 때 자동 백업. 무료 5GB는 대부분 부족, 6위안에 50GB로 업그레이드 권장.
  \item \textbf{Android}: 각 브랜드 클라우드 모두 백업 기능. 화웨이는 설정에서 "클라우드 백업" 검색, 샤오미는 "샤오미 클라우드 서비스", OPPO/vivo도 유사. 또 Google 자체 백업(설정 시스템 백업), 이것도 켜기.
  \item 클라우드 백업의 가장 큰 장점은 "신경 안 써도 됨", 자동 처리. 가장 큰 단점은 복원 시 다운로드가 오래 걸리고, 무료 공간이 보통 부족.
\end{itemize}

\booksection{셋, 분류 백업 전략: 중요한 것은 별도로 한 부 더 저장}

전체 백업은 최후 수단이지만, 일부는 너무 중요해 별도로 한 부 더 저장할 가치가 있습니다:

\booksubsection{사진과 비디오(가장 소중한 디지털 자산)}

\begin{itemize}
  \item \textbf{이중 보장 방안}: iCloud 사진/Google 포토 자동 동기화 켜기, 동시에 몇 달마다 수동으로 사진을 컴퓨터나 외장 하드에 복사. 이렇게 클라우드 1부, 로컬 1부, 상호 백업.
  \item \textbf{돈 안 쓰는 방안}: Google 포토는 무료 15GB, "압축 화질"로 설정해 공간 절약 가능. 또는 한 달마다 데이터 케이블로 사진을 전부 컴퓨터에 복사.
\end{itemize}

\booksubsection{연락처(잘 무시하지만 극히 중요)}

\begin{itemize}
  \item iPhone 연락처는 기본적으로 iCloud와 동기화, 동기화 켜면 추가 백업 불필요. 하지만 더 안전하게 하려면, iCloud.com에서 vCard 파일을 내보내기 컴퓨터에 저장.
  \item Android 연락처는 보통 Google 또는 브랜드 클라우드와 동기화. 연락처 앱 설정에서 .vcf 파일 내보내기, 완전한 백업 1부.
\end{itemize}

\booksubsection{위챗 채팅 기록(별도 처리, 최우선)}

이것은 반드시 따로 이야기해야 합니다, 전체 백업에 위챗 채팅 기록은 대부분 완전하지 않고, 기기 변경 도구로는 옮기지 못하기 때문.

\begin{itemize}
  \item \textbf{방법 1: 새 스마트폰으로 이전}. 위챗 나 설정 일반 채팅 기록 이전과 백업 이전 스마트폰/태블릿으로 이전. 새/구 스마트폰 같은 Wi-Fi 연결, QR 코드 스캔으로 전송. 주의: 이전 완료 전, 구형 스마트폰에서 위챗 로그아웃 금지.
  \item \textbf{방법 2: 컴퓨터로 백업}. 컴퓨터에 위챗 PC 버전 설치, 좌측 하단 세 줄 이전과 백업 백업과 복원 채팅 기록 컴퓨터로 백업. 이후 기기 변경 시 새 스마트폰으로 복원 가능.
  \item \textbf{방법 3: 클라우드로 백업}. 위챗도 iCloud 또는 스마트폰 클라우드 공간 백업 지원, 하지만 주의: 이 백업은 현재 스마트폰에서만 복원 가능, 크로스 디바이스 복원 불가. 기기 변경 시, 앞의 두 방법 사용해야.
\end{itemize}

\booksection{넷, 기기 변경 가이드: 구형 스마트폰에서 신형으로}

새 스마트폰으로 바꾸는 건 기쁜 일이지만, 데이터를 옮기지 못하면 두통으로 변합니다. 세 가지 상황별 대응:

\booksubsection{같은 브랜드 업그레이드(가장 간단)}

\begin{itemize}
  \item \textbf{iPhone 새 iPhone}: 부팅 활성화 시 구형 스마트폰을 옆에 두면, 자동으로 "빠른 시작" 팝업, 구형 스마트폰으로 신형의 패턴 스캔. 모든 데이터, 설정, 앱 레이아웃이 그대로 옮겨짐, Wi-Fi 비밀번호까지 자동 동반. 약 20-30분 소요.
  \item \textbf{Android 같은 브랜드}: 브랜드 자체 기기 변경 도구 사용. 화웨이는 "스마트폰 클론", 샤오미는 "샤오미 기기 변경", OPPO는 "스마트폰 이사", vivo는 "상호 전송". 단계 비슷: 새/구 스마트폰 같은 앱 설치, QR 코드 연결, 옮길 항목 선택, 전송 시작.
\end{itemize}

\booksubsection{크로스 브랜드 기기 변경(Android에서 Android, 다른 브랜드)}

서드파티 기기 변경 도구 또는 구 브랜드 자체 기기 변경 앱 사용(대부분 브랜드의 기기 변경 앱은 다른 Android 스마트폰으로 변경 지원). 옮길 수 있는 것: 연락처, 사진, 문자, 통화 기록, 일부 앱. 옮길 수 없는 것: 앱 내 로그인 상태(재로그인 필요), 위챗 채팅 기록(위챗 자체 이전 기능 사용).

\booksubsection{크로스 플랫폼 기기 변경(Android에서 iPhone, 또는 그 반대)}

\begin{itemize}
  \item \textbf{Android iPhone}: Apple의 "iOS로 전환" 앱, Android 스마트폰에서 다운로드 후, 새 iPhone 활성화 시 "Android 기기에서 데이터 전송" 선택. 연락처, 사진, 문자 등 기본 데이터 옮기기 가능, 앱 데이터는 옮길 수 없음. 위챗 채팅 기록은 위챗 자체 이전 기능 사용해야.
  \item \textbf{iPhone Android}: 각 Android 브랜드에 해당 이전 도구, 솔직히, 크로스 플랫폼 이전 경험은 같은 플랫폼만 못함. 가장 중요한 것부터 먼저 단독 처리: 연락처 내보내기, 사진 백업, 위챗 채팅 기록 이전. 나머지는 천천히.
\end{itemize}

\booksection{다섯, 백업의 황금 법칙: 3-2-1 원칙}

이것은 데이터 보안 분야의 고전 원칙, 단순하고 기억하기 쉽습니다:

\begin{itemize}
  \item \textbf{3}: 중요 데이터는 최소 3부 보존(원본 1부, 백업 2부).
  \item \textbf{2}: 백업은 두 가지 다른 저장 매체 사용(예: 클라우드 1부, 외장 하드 1부).
  \item \textbf{1}: 최소 1부 백업은 다른 물리적 위치에 보관(화재, 도난 등이 모든 데이터를 동시에 파괴하는 것 방지).
\end{itemize}

일반인에게 “클라우드 1부, 컴퓨터 1부”만 해도 이미 잘 한 겁니다. 너무 복잡하게 할 필요 없고, 핵심은 정말로 실행했다는 것.

\begin{tcbwarn}{위챗 채팅 기록 특별 주의}
위챗 채팅 기록은 어떤 스마트폰 기기 변경 도구로도 이전할 수 없습니다! 전체 백업 범위에 포함되지 않습니다. 반드시 위챗 자체의 “채팅 기록 이전” 기능 또는 컴퓨터 버전 위챗 백업 사용. 많은 사람이 새 스마트폰으로 바꾼 후 채팅 기록이 전부 없어진 것을 발견하는데, 그때는 이미 늦었습니다——구형 스마트폰이 이미 공장 초기화되었을 수 있으니까요.

\end{tcbwarn}

\begin{tcbpractice}{연습: 완전한 백업 훈련 한 번}
\begin{enumerate}
  \item iCloud 클라우드 백업 또는 브랜드 클라우드 백업 활성화 여부 확인, 마지막 백업 시간 확인.
  \item 수동 사진 백업: 데이터 케이블로 컴퓨터 연결, DCIM과 Pictures 폴더를 컴퓨터에 복사.
  \item 연락처 백업: iPhone 사용자는 iCloud.com에서 vCard 내보내기; Android 사용자는 연락처 앱 설정에서 .vcf 파일 내보내기.
  \item 위챗 채팅 기록 백업: 위챗 열기 나 설정 일반 채팅 기록 이전과 백업 컴퓨터로 백업.
  \item 내보낸 연락처 vcf와 사진 폴더를 외장 하드나 USB 메모리에 한 부 더 복사, "3-2-1"의 "한 부는 오프라인으로 다른 곳에" 완성.
\end{enumerate}

\end{tcbpractice}
